\input{preamble}

\title{Section 3.8 --- Rational Exponents --- Answer Key}

\begin{document}
\maketitle

As reference when completing this assignment, recall that we define
$$x^{p/q}=\sqrt[q]{x^p}=\Gp{\sqrt[q]{x}}^p$$

\section{Translating Radicals to Rational Exponents}
Rewrite each expression using rational exponents.  Use negative exponents rather than writing variable expressions in denominators.  Assume all variables represent positive real numbers.
\begin{smenum}
\mitemxxxx
{$25^{1/4}$}
{$r^{5/7}$}
{$6^{-2/7}$}
{$8n^{-9/7}$}
\mitemxxxx
{$5^{7/3}$}
{$m^{8/3}$}
{$(5w+3)^{5/7}$}
{$(2x^2-1)^{3/4}$}

\end{smenum}

\section{Translating Rational Exponents to Radicals}
Rewrite each expression in radical notation. (For number bases, don't evaluate, just rewrite.) Your answer should include only positive integer exponents.
\begin{smenum}
\mitemxxxx
{$\sqrt[3]{27}$}
{$\sqrt[4]{y}$}
{$-\sqrt[4]{x^3}$}
{$8\sqrt[3]{n^2}$}
\mitemxxxx
{$\sqrt{x^5}$}
{$\dfrac{3}{\sqrt[3]{y^2}}$}
{$\dfrac{1}{\sqrt{4x}}$}
{$\dfrac{1}{\sqrt[5]{(xy)^8}}$ or $\dfrac{1}{\sqrt[5]{x^8y^8}}$}
\end{smenum}

\section{Computing Rational Powers}
Compute each power by hand. You can use the following reference:
\begin{center}
\begin{tabularx}{0.75\linewidth}{CCCCCC}
\hline
$x$&$x^2$&$x^3$&$x^4$&$x^5$&$x^6$\\\hline
1&\ph1&\ph\ph1&\phc\ph\ph1&\phc\ph\ph1&\ph\phc\ph\ph1\\
2&\ph4&\ph\ph8&\phc\ph16&\phc\ph32&\ph\phc\ph64\\
3&\ph9&\ph27&\phc\ph81&\phc243&\ph\phc729\\
4&16&\ph64&\phc256&1,024&\ph4,096\\
5&25&125&\phc625&3,125&15,625\\
6&36&216&1,296&7,776&46,656\\\hline
\end{tabularx}
\end{center}
If the given power is not a real number, write ``undefined.'' All answers for this section are rational.
\begin{smenum}
\mitemxxxx
{$64^{1/6}=2$}
{$\Gp{\dfrac{1}{8}}^{-2/3}=4$}
{$(-32)^{-3/5}=-\dfrac{1}{8}$}
{Undefined}
\mitemxxxx
{$\dfrac{3}{2}$}
{Undefined}
{$\dfrac{27}{125}$}
{$\dfrac{625}{1296}$}
\end{smenum}
Write each expression in simplified radical form, including rationalizing any denominators.
\begin{smenum}
\mitemxxxx
{$\sqrt[3]{4}$}
{$8\sqrt[3]{2}$}
{$\frac{\sqrt{3}}{9}$}
{$\frac{\sqrt{3}}{18}$}
\end{smenum}

\section{Simplifying with Rational Exponents}
Simplify each expression as much as possible. Your answer should include a single occurrence of each variable and only positive exponents. Leave your answer in terms of rational exponents. To avoid concerns about where expressions are defined, assume all variables represent positive real numbers.
\begin{smenum}
\mitemxxxx
{$x^{1/3}y^{1/3}z^{2/3}$}
{$x^{1/4}$}
{$x^{4/3}$}
{$\dfrac{y^{3/2}}{x^{14/5}}$}
\mitemxxxx
{$\dfrac{y^{1/12}}{x^{1/6}}$}
{$\dfrac{y^{1/6}}{x^{5/6}z^{2/3}}$}
{$\dfrac{1}{x^{2/3}y^{12/5}}$}
{$\dfrac{y^{3/14}z^{5/14}}{x^{4/21}}$}
\end{smenum}

\end{document}
